\subsection{LQR with integral action}

The final LQ controller used in simulation and on the real setup includes integral action on the
position error. The simple LQR without integrator is not retained, because reference tracking would
depend on an additional static scaling and would be more sensitive to modelling errors and steady-state
offsets. For this reason, the design is performed directly on the augmented model.

The integral state is defined as
\begin{equation}
    \dot{\xi} = \Delta x_{ref} - \Delta x.
\end{equation}
The augmented state is
\begin{equation}
    \vect{x}_{aug} =
    \begin{bmatrix}
        \Delta x & \Delta \dot{x} & \Delta i & \xi
    \end{bmatrix}^T.
\end{equation}
The augmented matrices are
\begin{equation}
    A_{aug} =
    \begin{bmatrix}
        A_{ext} & 0 \\
        -C_{pos} & 0
    \end{bmatrix},
    \qquad
    B_{aug} =
    \begin{bmatrix}
        B_{ext} \\
        0
    \end{bmatrix}.
\end{equation}

The weights used for the nominal simulation design in \texttt{SS\_model\_LQ.m} are
\begin{equation}
    Q_{aug} = \mathrm{diag}(5000, 5\times10^{-1}, 10^{-9}, 9\times10^{4}),
    \qquad
    R_{aug} = 5\times10^{-4}.
\end{equation}
The resulting gain is
\begin{equation}
    K_{aug}
    =
    \mathrm{lqr}(A_{aug}, B_{aug}, Q_{aug}, R_{aug})
    =
    \begin{bmatrix}
        K_{x,lqi} & K_{i,lqi}
    \end{bmatrix}.
\end{equation}
The control law is therefore
\begin{equation}
    \Delta i_{ref}
    =
    -K_{x,lqi}
    \begin{bmatrix}
        \Delta x & \Delta \dot{x} & \Delta i
    \end{bmatrix}^T
    -K_{i,lqi}\xi.
\end{equation}
